\begin{landscape}
\begingroup
\small
\singlespacing
\setlength{\parindent}{0pt}
\setlength{\tabcolsep}{5pt}
\renewcommand{\arraystretch}{1.12}
\centering
\captionof{table}{Les hypothèses diagnostiques}
\label{tab:hypothesesDiagnostiques}
\noindent\begin{tabularx}{\linewidth}{|L{2.8cm}|L{4.5cm}|L{4.9cm}|L{6.4cm}|Y|}
  \hline
  Diagnostic & Définition & Symptômes & Signes cliniques & Remarques \\
  \hhline{|=====|}
  Kyste hémorragique du corps jaune &
  Kyste fonctionnel, dépendant du cycle menstruel. &
  Dysménorrhée avec douleur unilatérale lors de la phase lutéale. &
  \begin{tableitemize}
    \item Palpation abdominale : sensible ;
    \item Toucher vaginale : masse annexielle palpée, mobilisation utérine douloureuse.
  \end{tableitemize} &
  Rupture du kyste $\rightarrow$~défense abdominale, signe du cri du Douglas. \\
  \hline
  Endométriome &
  Endomètre ectopique au niveau ovarien. &
  Dysménorrhée, douleurs pelviennes chroniques, dyspareunie profonde, troubles de la fertilité. &
  \begin{tableitemize}
    \item Palpation abdominale : sensible ;
    \item Toucher vaginale : masse annexielle palpable, mobilisation utérine douloureuse (adhérences suspectées), cul-de-sac du Douglas sensible.
  \end{tableitemize} &
  Les symptômes associés à d’autres systèmes sont observés (nycturie, dyschésie, trouble du transit, \ldots). Cependant, aucune explication scientifique ne parvient à établir un lien de causalité entre les symptômes et le kyste. \\
  \hline
  Kyste dermoïde &
  Tératome, tumeur bénigne d'origine embryonnaire, fréquente chez les femmes entre 16 et 55 ans. &
  Asymptomatique ou légère douleur pelvienne, métrorragie, dyspareunie et symptômes liés à une compression éventuelle. &
  \begin{tableitemize}
    \item Palpation abdominale : normal ;
    \item Toucher vaginale : masse annexielle peut-être palpable.
  \end{tableitemize} &
  Torsion ovarienne = urgence gynécologique $\rightarrow$~douleur abdominale aiguë et intense. \\
  \hline
  Kyste para-ovarien &
  Kyste bénin avec pour origine le mésosaplinx. &
  Asymptomatique ou symptômes liés à une compression éventuelle. &
  \begin{tableitemize}
    \item Palpation abdominale : normal ;
    \item Toucher vaginale : masse annexielle peut-être palpable.
  \end{tableitemize} &
  \\
  \hline
  Tumeur ovarienne &
  Cancer ovarien malin. Facteurs favorisants : G0P0, ménarche précoce, cycle menstruel spontané, \ldots &
  Stade précoce : paucisymptomatique.\newline
  Stade avancé : AEG, perte de poids, fatigue, \ldots &
  \begin{tableitemize}
    \item Palpation abdominale : signes de carcinose (ventre de batracien, \ldots);
    \item toucher vaginale : masse annexielle palpable voire d'autres masses selon le stade d'extension.
  \end{tableitemize} &
  Ce diagnostic doit être exclu vu son pronostic dramatique. \\
  \hline
\end{tabularx}
\par
\endgroup
\end{landscape}
